% Part: first-order-logic
% Chapter: sequent-calculus
% Section: soundness-identity

\documentclass[../../../include/open-logic-section]{subfiles}

\begin{document}

\olfileid[fr]{fol}{seq}{sid}

\olsection{Correction en présence de l'\usetoken{s}{identity}}

\begin{prop}
Le calcul $\Log{LK}$ muni des séquents initiaux et des règles pour
l'égalité est correct.
\end{prop}

\begin{proof}
Les séquents initiaux de la forme ${} \Sequent \eq[t][t]$ sont
valides : pour toute !!{structure}~$\Struct M$, on a
$\Sat{M}{\eq[t][t]}$. Le terme $t$ est supposé clos, c'est-à-dire
sans variable ; les valuations des variables n'interviennent donc pas.

Supposons que la dernière inférence d'une !!{derivation} soit la
règle $=$ dont la prémisse est
$\eq[t_1][t_2], \Gamma \Sequent \Delta, !A(t_1)$ et la conclusion
$\eq[t_1][t_2], \Gamma \Sequent \Delta, !A(t_2)$.
Considérons une !!{structure}~$\Struct M$.
Pour établir la validité de la conclusion, il faut montrer que,
si $\Sat{M}{\eq[t_1][t_2]}$ et $\Sat{M}{\Gamma}$, alors
$\Sat{M}{!C}$ pour un certain $!C \in \Delta$ ou
$\Sat{M}{!A(t_2)}$.

Par hypothèse de récurrence, la prémisse est valide.
Si $\Sat{M}{\eq[t_1][t_2]}$ et $\Sat{M}{\Gamma}$, on a donc
soit (a) $\Sat{M}{!C}$ pour un certain $!C \in \Delta$,
soit (b) $\Sat{M}{!A(t_1)}$.
Le cas (a) donne le résultat voulu. Dans le cas (b), soit $s$
une valuation des variables telle que $s(x) = \Value{t_1}{M}$.
D'après la \olref[syn][ass]{prop:sentence-sat-true},
$\Sat{M}{!A(t_1)}[s]$.
Comme $\varAssign{s}{s}{x}$, la
\olref[syn][ext]{prop:ext-formulas} donne $\Sat{M}{!A(x)}[s]$.
Puisque $\Sat{M}{\eq[t_1][t_2]}$, on a
$\Value{t_1}{M} = \Value{t_2}{M}$, et donc
$s(x) = \Value{t_2}{M}$.
En appliquant de nouveau la \olref[syn][ext]{prop:ext-formulas},
on obtient $\Sat{M}{!A(t_2)}[s]$.
Enfin, la \olref[syn][ass]{prop:sentence-sat-true} entraîne
$\Sat{M}{!A(t_2)}$.
L'autre règle pour l'égalité se traite de la même façon, en
intervertissant les rôles de $t_1$ et de $t_2$ dans cet argument :
l'égalité de leurs valeurs est symétrique.\footnote{Note de l'édition
française : cette dernière phrase explicite le cas de la seconde
orientation de la règle, que la source ne développe pas.}
\end{proof}

\end{document}
